\documentclass[10pt]{article}
\usepackage[T1]{fontenc}
\usepackage[utf8]{inputenc}
\usepackage[french]{babel}
\usepackage{amsmath,amssymb,mathrsfs,mathpazo}
\usepackage{geometry}
\geometry{a4paper,margin=2cm}
\usepackage{graphicx,array,tabularx,tabularray,colortbl,multirow,multicol,booktabs,enumitem,pifont,ulem,fancybox,diagbox,cancel}
\usepackage{pstricks,pstricks-add,pst-plot,pst-tree,pst-node,pst-text,pst-func,pst-math,pst-bspline,pst-3dplot,pst-eucl,pst-solides3d}
\let\pscalculate\undefined
\usepackage{pst-bezier}
\usepackage{tikz,pgfplots,tkz-tab}
\usetikzlibrary{arrows,arrows.meta,patterns,calc,positioning,trees,decorations.markings,shapes.geometric,shapes.misc,shapes.arrows,intersections,decorations.pathmorphing,decorations.pathreplacing,backgrounds,shadows,babel}
\pgfplotsset{compat=1.18}
\usepgfplotslibrary{fillbetween}
\usepackage[np]{numprint}
\usepackage{listings,makecell,pas-tableur,twoopt}
\let\barre\undefined
\usepackage[pstricks]{tablvar}
\definecolor{darkgreen}{RGB}{0,100,0}
\providecommand{\red}{\color{red}}
\providecommand{\blue}{\color{blue}}
\providecommand{\green}{\color{green}}
\providecommand{\black}{\color{black}}
\providecommand{\white}{\color{white}}
\providecommand{\gray}{\color{gray}}
\providecommand{\cyan}{\color{cyan}}
\providecommand{\magenta}{\color{magenta}}
\providecommand{\yellow}{\color{yellow}}
\providecommand{\violet}{\color{violet}}
\providecommand{\orange}{\color{orange}}
\providecommand{\euro}{\texteuro}
\providecommand{\EUR}{\texteuro}
\providecommand{\ieme}{\textsuperscript{e}}
\providecommand{\degres}{\ensuremath{^\circ}}
\providecommand{\degree}{\ensuremath{^\circ}}
\providecommand{\vv}[1]{\overrightarrow{#1}}
\providecommand{\Coord}[1]{#1}
\providecommand{\decosix}{}
\providecommand{\starredbullet}{\ensuremath{\bullet}}
\pagestyle{empty}
\setlength{\parindent}{0pt}
\usepackage{adjustbox}\DeclareUnicodeCharacter{2236}{:}\DeclareUnicodeCharacter{2190}{\ensuremath{\leftarrow}}
\providecommand{\R}{\mathbb{R}}
\providecommand{\N}{\mathbb{N}}
\providecommand{\Z}{\mathbb{Z}}
\providecommand{\Q}{\mathbb{Q}}
\providecommand{\D}{\mathbb{D}}
\providecommand{\C}{\mathbb{C}}
\providecommand{\e}{{\mathrm{e}}}
\providecommand{\vect}[1]{\overrightarrow{#1}}
\providecommand{\vectt}[1]{\overrightarrow{\mathrm{#1}}}
\providecommand{\vvt}[1]{\overrightarrow{\mathrm{#1}}}
\providecommand{\np}[2][]{{#2~#1}}
\providecommand{\vv}[1]{{\overrightarrow{#1}}}
\renewcommand{\d}{{\mathrm{d}}}
\providecommand{\degres}{{^\circ}}
\providecommand{\degre}{{^\circ}}
\providecommand{\blue}{\color{blue}}
\providecommand{\red}{\color{red}}
\providecommand{\enskip}{\quad}
\providecommand{\llbracket}{\lbrack\!\lbrack}
\providecommand{\rrbracket}{\rbrack\!\rbrack}
\providecommand{\euro}{{\text{€}}}
\providecommand{\no}{{\text{n°}}}
\providecommand{\interval}[2]{{[#1;#2]}}
\providecommand{\intervFO}[2]{{[#1;#2[}}
\providecommand{\dotfill}{\ldots}

\providecommand{\ds}{\displaystyle}
\providecommand{\textding}[1]{#1}
\providecommand{\Oij}{\(\mathrm{O};\vec{\imath},\vec{\jmath}\)}
\providecommand{\Oijk}{\(\mathrm{O};\vec{\imath},\vec{\jmath},\vec{k}\)}
\providecommand{\Ouv}{\(\mathrm{O};\vec{u},\vec{v}\)}
\providecommand{\intervFF}[2]{\ensuremath{[#1;#2]}}
\renewcommand{\euro}{\ifmmode\text{\texteuro}\else\texteuro\fi}
\DeclareUnicodeCharacter{00B0}{\ifmmode^{\circ}\else\textdegree\fi}
\DeclareUnicodeCharacter{20AC}{\ifmmode\text{\texteuro}\else\texteuro\fi}
% Énoncés : Éducation nationale. Transcription et corrigés : APMEP, Denis Vergès et contributeurs. Découpage et classement : CoopMaths, Sébastien Lozano, Éric Elter et contributeurs. Adaptation de la présentation : Maths Explorer.
% CC BY-SA : https://coopmaths.fr/www/about/
% https://coopmaths.fr/alea/static/bac/2026/tex/bac_2026_06_sujet1_asie_4.tex
\begin{document}
\noindent\textit{APMEP / CoopMaths. CC BY-SA. Présentation adaptée par Maths Explorer.}\par\medskip
\begin{enumerate}
\item On considère la fonction $g$ définie sur l'intervalle $[0~;~2\pi]$ par :
\[
g(x)=x\cos(x)-\sin(x)
\]

On admet que la fonction $g$ est dérivable sur l'intervalle $[0~;~2\pi]$ et on note $g'$ sa dérivée.

\begin{enumerate}
\item Montrer que pour tout réel $x$ de l'intervalle $[0~;~2\pi]$, on a :
\[
g'(x)=-x\sin(x)
\]

\item On donne le tableau de variations de la fonction $g$ sur l'intervalle $[0~;~2\pi]$ ci-dessous.

Justifier chacun des éléments qui figurent dans ce tableau de variations.

\begin{center}
\begin{tikzpicture}[x=1cm,y=1cm]
\draw (0,0) rectangle (8,2.4);
\draw (1.2,0) -- (1.2,2.4);
\draw (0,1.6) -- (8,1.6);
\node at (0.6,2) {$x$};
\node at (2,2) {$0$};
\node at (4.6,2) {$\pi$};
\node at (7.2,2) {$2\pi$};
\node at (0.6,0.8) {$g$};
\node at (2,1.25) {$0$};
\node at (4.6,0.35) {$-\pi$};
\node at (7.2,1.25) {$2\pi$};
\draw[->] (2.25,1.05) -- (4.25,0.55);
\draw[->] (4.95,0.55) -- (6.95,1.05);
\end{tikzpicture}
\end{center}

\item Montrer qu'il existe une unique valeur réelle $\alpha$ dans l'intervalle $[\pi~;~2\pi]$ telle que $g(\alpha)=0$.

\item En déduire le tableau de signes de la fonction $g$ sur l'intervalle $[0~;~2\pi]$.
\end{enumerate}

\item On considère la fonction $f$ définie sur l'intervalle $]0~;~2\pi]$ par :
\[
f(x)=\dfrac{\sin(x)}{x}
\]

On admet que la fonction $f$ est dérivable sur l'intervalle $]0~;~2\pi]$ et on note $f'$ sa dérivée.

\begin{enumerate}
\item Montrer que pour tout réel $x$ de l'intervalle $]0~;~2\pi]$ on a :
\[
f'(x)=\dfrac{g(x)}{x^2}
\]

\item Étudier le signe de la fonction $f'$ sur l'intervalle $]0~;~2\pi]$.

\item En déduire le sens de variation de la fonction $f$ sur l'intervalle $]0~;~2\pi]$.

\item Déterminer la limite de $f$ en $0$. On pourra utiliser le taux d'accroissement de la fonction sinus en $0$.
\end{enumerate}

\item On considère deux nombres réels $r$ et $s$ qui vérifient l'inégalité : $0 < r < s < \pi$.

Montrer que :
\[
\dfrac{r}{s}<\dfrac{\sin(r)}{\sin(s)}
\]
\end{enumerate}

\end{document}
